\chapter{Transformer autoregressivi}
\label{cap:transformer}

Questo capitolo richiama l'architettura Transformer nella variante decoder-only
usata dai modelli linguistici autoregressivi \parencite{vaswani2017attention},
limitandosi a quanto serve per localizzare il punto in cui l'esperimento~01
sostituisce uno strato lineare con un circuito. La trattazione segue
l'implementazione del repository e ne segnala le divergenze dalla formulazione
canonica.

L'implementazione classica di partenza riprende la struttura didattica costruita
da Karpathy nel tutorial \emph{Let's build GPT: from scratch}, parte del corso
\emph{Neural Networks: Zero to Hero} \parencite{karpathyZeroToHero}. Il
repository nanoGPT dello stesso autore ha costituito un ulteriore riferimento
progettuale \parencite{karpathyNanoGPT}; il modello impiegato qui non ne è però
una replica, perché conserva teste di attenzione separate, tokenizzazione a
caratteri e dimensioni molto più ridotte.

\paragraph{Convenzione di notazione.} $\dbatch$ è la dimensione di batch, $\dseq$
la lunghezza di sequenza (\code{block\_size}), $\demb$ la dimensione di embedding
(\code{n\_embd}), $\nhead$ il numero di teste, $\dhead = \demb/\nhead$ la
dimensione per testa (\code{head\_size}), $\dvoc$ la cardinalità del vocabolario.
La divisibilità $\demb \bmod \nhead = 0$ è verificata all'istanziazione della
configurazione.

\section{Embedding e tokenizzazione}
\label{ssec:tokenizer}

L'ingresso dei blocchi è la somma di due tabelle apprese, quella dei token e
quella delle posizioni:
\begin{equation}
  x \;=\; E[\text{idx}] + P[0{:}\dseq],
  \qquad E \in \R^{\dvoc \times \demb},\quad
  x \in \R^{\shape{\dbatch, \dseq, \demb}} .
  \label{eq:input-gpt}
\end{equation}
Le proiezioni e il calcolo dei punteggi non rappresentano da soli la posizione
dei token, che viene quindi aggiunta esplicitamente. Nel codice entrambe le
tabelle sono
\code{nn.Embedding}, quindi l'embedding posizionale è appreso anziché sinusoidale;
ha esattamente \code{block\_size} righe, per cui il modello non estrapola oltre la
lunghezza di contesto con cui è stato addestrato e in generazione la finestra
viene troncata.

Nelle run definitive si usa \classe{CharTokenizer}, che costruisce il vocabolario
dai caratteri distinti del corpus. Il repository include anche
\classe{BiCharTokenizer}, usato soltanto nei pilot storici. Nessuno dei due è un
tokenizer a sottoparole addestrato.

\begin{limite}
La perplexity calcolata su un vocabolario di caratteri non è confrontabile con
quella di un modello a sottoparole, perché cambia l'unità su cui la
verosimiglianza è normalizzata. I valori riportati nel \cref{cap:risultati} vanno
letti \emph{soltanto} nel confronto reciproco fra la variante classica e quella
quantistica addestrate sullo stesso corpus e con lo stesso tokenizer.
\end{limite}

\section{Self-attention e teste multiple}
\label{sec:self-attention}

Dato $x \in \R^{\shape{\dbatch, \dseq, \demb}}$ e tre matrici di proiezione
$W_Q, W_K, W_V \in \R^{\demb \times \dhead}$, si pone $Q = xW_Q$, $K = xW_K$,
$V = xW_V$, e l'uscita di una testa è $\softmax(S)\,V$ con
\begin{equation}
  S \;=\; \frac{QK^{\top}}{\sqrt{\dhead}} ,
  \qquad S \in \R^{\shape{\dbatch, \dseq, \dseq}} .
  \label{eq:scores}
\end{equation}
Se le componenti di $q$ e $k$ sono indipendenti con media nulla e varianza
unitaria, il prodotto scalare di due vettori $d$-dimensionali ha varianza $d$.
La divisione per $\sqrt{\dhead}$ evita quindi che l'aumento della dimensione
spinga la $\softmax$ verso una distribuzione quasi degenere. Nel
codice la scala è \code{self.scale = head\_size**-0.5}, uguale per entrambe le
varianti. Il pilot iniziale usava per errore $1/\sqrt{\demb}$; quella run è stata
esclusa dalle run definitive e la correzione applicata prima di fissare il preset
\code{thesis}.

L'obiettivo autoregressivo richiede poi che la predizione alla posizione $t$
dipenda solo dalle posizioni $t' \le t$. La maschera causale pone $S_{t,t'} =
-\infty$ per $t' > t$ \emph{prima} della $\softmax$, così che $\exp(-\infty) = 0$
produca peso esattamente nullo dopo la normalizzazione. La matrice triangolare è
registrata con \code{register\_buffer}: viene salvata nel checkpoint e spostata
col modello fra CPU e GPU, ma non è un parametro e non entra nel conteggio del
\cref{cap:risultati}.

L'attenzione multi-testa esegue $\nhead$ attenzioni parametrizzate separatamente
in sottospazi di dimensione $\dhead$, ne concatena le uscite e applica una
proiezione finale $W_O \in \R^{\demb \times \demb}$ seguita da dropout.
Nell'implementazione le
teste sono moduli separati in un \code{nn.ModuleList} di istanze di \classe{Head},
e non sono fuse in un'unica moltiplicazione di matrici come nelle implementazioni
ottimizzate. La variante quantistica richiede infatti un modulo distinto per ogni
proiezione di ogni testa, con circuito e pesi propri. La stessa struttura è usata
nel modello classico, così che la differenza architetturale riguardi il tipo di
proiezione.

Nell'esperimento~01 la dimensione della testa coincide con il numero di qubit. La
proprietà \code{GPTConfig.n\_qubits} è definita come
\begin{equation}
  \nq \;=\; \frac{\demb}{\nhead} \;=\; \dhead ,
  \label{eq:nqubits-gpt}
\end{equation}
Il metodo \code{Head.\_get\_layer} verifica questa uguaglianza e segnala
dimensioni incompatibili. Il preset definitivo \code{thesis} usa $\demb=8$ e
$\nhead=2$, quindi $\nq=4$ qubit per proiezione.

\section{Rete feed-forward}
\label{sec:ffn}

Ogni blocco contiene una rete densa applicata indipendentemente a ciascuna
posizione, $\text{FFN}(x) = \text{Dropout}\big(W_2\,\mathrm{ReLU}(W_1 x + b_1) +
b_2\big)$ con espansione di fattore quattro, $W_1 \in \R^{\demb \times 4\demb}$ e
$W_2 \in \R^{4\demb \times \demb}$, come in \textcite{vaswani2017attention}.
ReLU introduce la non-linearità nella FFN; il blocco comprende inoltre la
$\softmax$ dell'attenzione e le normalizzazioni.

Con il preset \code{thesis}, $\demb=8$, la FFN contiene \num{552} parametri per
blocco. Le proiezioni $Q,K,V$ classiche ne contengono complessivamente \num{192},
la proiezione di uscita 72 e le due \code{LayerNorm} 32. La FFN rappresenta
quindi circa il \SI{65,1}{\percent} degli 848 parametri del blocco e resta
classica in entrambe le varianti.

Il VQC modifica soltanto le proiezioni $Q,K,V$. In assenza di differenze
misurabili, i dati non permettono di distinguere tra inefficacia della
sostituzione ed effetto attenuato dal resto dell'architettura
(\cref{sec:capacita}).

\section{Connessioni residue e normalizzazione}
\label{sec:residui}

Ogni sottostrato è applicato in forma additiva, $x \leftarrow x + f(x)$, e offre
al gradiente un percorso diretto attraverso il ramo identità. Il blocco adotta la
variante \emph{pre-norm}, \code{x = x + self.sa(self.ln1(x))} seguito da
\code{x = x + self.ffwd(self.ln2(x))}, con una \code{LayerNorm} finale prima della
testa di linguaggio: rispetto al \emph{post-norm} originale, questa forma lascia
il ramo identità privo di normalizzazioni. L'addestramento usa un learning rate
costante, senza \emph{warmup} né scheduler.

Il ramo identità non è limitato, mentre l'uscita del VQC è confinata in
$[-1,1]^{\nq}$ dai valori di aspettazione (\cref{sec:misurazione}). Se la norma
della rappresentazione residua cresce con la profondità, l'uscita del circuito
può quindi diventare piccola rispetto al ramo identità. I dati raccolti non
misurano direttamente questo rapporto; la \cref{sec:saturazione} lo considera
come possibile limite architetturale.

\section{Obiettivo e metrica}

La verosimiglianza di una sequenza è fattorizzata come
$p(x_{1:\dseq}) = \prod_{t} p(x_t \mid x_{<t})$, e il modello minimizza la
cross-entropy media per token $\mathcal{L}$. I logits di forma
$\shape{\dbatch, \dseq, \dvoc}$ sono rimodellati in
$\shape{\dbatch\cdot\dseq, \dvoc}$ e confrontati con i target appiattiti tramite
\code{F.cross\_entropy}. La perplexity è $\PPL = \exp(\mathcal{L})$, e si
interpreta come il numero effettivo di alternative equiprobabili fra cui il
modello esita a ogni token.

La loss di validation è prodotta da \code{estimate\_loss}, che media su
\code{eval\_iters} batch estratti a caso sotto \code{torch.no\_grad()}. Le
varianti appaiate sono valutate sugli stessi batch; la dispersione fra seed
riportata nel \cref{cap:risultati} comprende invece l'intera variabilità delle
esecuzioni. Il metodo \code{QuantumGPT.generate} campiona infine in modo
iterativo da \code{softmax(logits[:, -1, :])} con \code{torch.multinomial}, senza
\emph{top-$k$} né temperatura.

\section{Sostituzione delle proiezioni}
\label{sec:corrispondenza-gpt}
\label{ssec:punto-sostituzione}

La differenza architetturale riguarda la scelta, per ciascuna delle proiezioni
\code{key}, \code{query} e \code{value} di ogni testa, fra
\classe{QuantumLayerAdapter} e un ordinario \code{nn.Linear}, operata da
\code{Head.\_get\_layer} in
\file{01\_quantum\_gpt/src/model.py}. Nel percorso quantistico ciascuna proiezione
segue la catena \code{Linear} $\to\tanh\cdot\pi\to$ \code{AngleEmbedding} $\to$
\code{StronglyEntanglingLayers} $\to\expval{Z}$.

Restano classici il prodotto $QK^{\top}$, il fattore di scala, la maschera
causale, la $\softmax$, l'aggregazione dei valori, la proiezione di uscita, la
rete feed-forward, gli embedding e la testa di linguaggio. Nel seguito il modello
è quindi descritto come Transformer con proiezioni $Q,K,V$ quantistiche.

\paragraph{Implementazioni correlate.} Quixer combina primitive come Linear
Combination of Unitaries e Quantum Singular Value Transform in un modello di
linguaggio quantistico \parencite{khatri2024quixer}, e il codice di simulazione
con baseline Transformer, LSTM e FNet è disponibile nel repository ufficiale
\parencite{quantinuumQuixerRepo}. Il repository \code{rdisipio/qtransformer}
fornisce invece un esempio indipendente di Transformer ibrido
\parencite{disipioQTransformerRepo}. Sono riferimenti di confronto e non
dipendenze del presente codice.

\begin{table}[htbp]
\centering
\caption[Parametri per proiezione $Q$, $K$ o $V$]{%
Parametri addestrabili di una singola proiezione, preset \code{thesis}
($\demb = 8$, $\nhead = 2$, $\nqlayers = 2$). La variante quantistica ha più
parametri addestrabili della baseline classica.}
\label{tab:parametri-proiezione}
\small
\begin{tabular}{@{}l l r@{}}
\toprule
\textbf{Variante} & \textbf{Composizione} & \textbf{Parametri} \\
\midrule
Classica    & $\demb \times \dhead$, senza bias & 32 \\
\midrule
Quantistica & adapter: $\demb \times \nq + \nq$ & 36 \\
            & VQC: $\nqlayers \cdot \nq \cdot 3$ & 24 \\
            & \textbf{totale} & \textbf{60} \\
\bottomrule
\end{tabular}
\end{table}

Sul modello completo i totali sono \num{2945} parametri classici e \num{3281}
quantistici, cioè \num{336} in più, pari all'$11{,}4\,\%$. La corrispondenza fra
gli elementi teorici del capitolo e le classi che li implementano è
nell'appendice~\ref{app:circuiti}.
